\begin{convention}[Index suppression and aggregation from the root variable]
\label{conv:notation}
Every variable in the model is defined with an ordered list of index--set
pairs $(i_1 \in I_1,\, i_2 \in I_2,\, \ldots,\, i_n \in I_n)$, called its
\emph{root definition}; the variable together with its root definition is the
\emph{root variable}. All variables in Table~\ref{tab:nomenclature} are root
variables.

\medskip
\noindent\textbf{Index suppression.}
Writing a root variable $\phi$ with a strict subset of its indices, preserving
their original order, is called \emph{index suppression}. The indices that
appear are the \emph{retained} indices; those omitted are the \emph{suppressed}
indices. Because the root definition is fixed, the set each suppressed index
ranges over is already known — it is the set paired with that index in the
root definition.

Index suppression defines an aggregation group, $\mathcal{G}$. To formalize this, let $K \subset \{1, \dots, n\}$ denote the index positions that are retained, and let $S = \{1, \dots, n\} \setminus R$ denote the index positions that are suppressed.

For a given tuple of fixed values for the retained indices, denoted as ${i_{r_1}\ldots,i_{r_k} $, the aggregation group is the set of all root variable instances formed by locking the retained indices to these fixed values, while the suppressed indices vary freely over their root-defined sets:$$\mathcal{G}_{(\bar{\imath}_r)_{r \in R}} \coloneqq \left\{ \phi_{i_1,\ldots,i_n} \;\middle|\; i_r = \bar{\imath}_r \text{ for all } r \in R, \text{ and } i_s \in I_s \text{ for all } s \in S \right\}$$
For each fixed tuple of retained index values, index suppression defines the
\emph{aggregation group} as the set of all root variable instances obtained
by varying the suppressed indices over their root-defined sets:
\begin{equation}
  \mathcal{G}_{i_{r_1},\ldots,i_{r_k}}
  \;\coloneqq\;
  \bigl\{\,
    \phi_{i_1,\ldots,i_n}
    \;:\;
    i_{r_q} \text{ fixed for each retained index } i_{r_q},\quad
    i_s \in I_s \text{ for each suppressed index } i_s
  \,\bigr\},
  \label{eq:group}
\end{equation}
where $i_{r_1},\ldots,i_{r_k}$ are the retained indices, fixed at given
values, and $i_{s_1},\ldots,i_{s_l}$ are the suppressed indices, varying
freely over their root-defined sets. An \emph{aggregation rule} then maps $\mathcal{G}_{i_{r_1},\ldots,i_{r_k}}$
to a single derived variable indexed only by the retained indices. The rules
are defined below. Figure~\ref{fig:index-suppression} illustrates the grouping.

When two indices share the same root set~$S$, suppressing one is ambiguous:
the position of the surviving index alone does not identify which of the two
was suppressed. A dot~$(\,\cdot\,)$ marks the suppressed position explicitly,
restoring the ordering information:
\[
  \phi_{\,\cdot\; j\,m} \quad \text{(first index of $S$ suppressed)},
  \qquad
  \phi_{i\;\cdot\; m} \quad \text{(second index of $S$ suppressed)}.
\]

\medskip
\noindent\textbf{Rule 1 (Sum aggregation).}
The sum aggregation rule maps $\mathcal{G}_{i_1,\ldots,i_k}$ to the sum of
all its elements, yielding the derived variable
\[
  \widetilde{\phi}_{i_1,\ldots,i_k}
  \;\coloneqq\;
  \sum_{j_1 \in J_1}\cdots\sum_{j_l \in J_l}
  \phi_{i_1,\ldots,i_k,\,j_1,\ldots,j_l}.
\]
For example, suppressing subsets of indices from $NF_{cijmpr}$ gives:
\begin{align*}
  \widetilde{NF}_{cijm}
    &\coloneqq \sum_{p \in P}\sum_{r \in R} NF_{cijmpr}, &
  \widetilde{NF}_{cij}
    &\coloneqq \sum_{m \in \mathcal{M}}\sum_{p \in P}\sum_{r \in R} NF_{cijmpr}, &
  \widetilde{NF}
    &\coloneqq \sum_{i \in V}\sum_{j \in V}
               \sum_{m \in \mathcal{M}}\sum_{p \in P}\sum_{r \in R} NF_{cijmpr}.
\end{align*}
For continuous and integer variables this is the natural interpretation. For
binary variables the result is an \emph{integer count} of how many elements
of $\mathcal{G}$ equal~1; it is not itself binary. For example,
$\widetilde{X}_{cij} \coloneqq \sum_{m \in \mathcal{M}} X_{cijm}$ counts
active modes for commodity type~$c$ on lane $(i,j)$ and takes values in
$\{0,1,\dots,|\mathcal{M}|\}$.

A derived variable may also be introduced explicitly into the model together
with its defining equality constraint; for instance:
\begin{equation}
  \widetilde{NF}_{cijm}
    \coloneqq \sum_{p \in P}\sum_{r \in R} NF_{cijmpr}.
  \label{eq:aux-sum}
\end{equation}

\medskip
\noindent\textbf{Rule 2 (OR-aggregated binary).}
The sum aggregate $\widetilde{X}_{cij}$ of a binary root variable is an
integer count over $\mathcal{G}_{cij}$ and therefore not itself binary. When
a proper binary is needed that equals~1 if and only if \emph{at least one}
element of $\mathcal{G}_{cij}$ equals~1, a new binary variable
$\widecheck{X}_{cij}$ is introduced and linked to the sum aggregate via:
\begin{align}
  \widecheck{X}_{cij} &\leq \widetilde{X}_{cij},
    \label{eq:or-m-lb} \\
  \widetilde{X}_{cij} &\leq |\mathcal{M}|\cdot\widecheck{X}_{cij}.
    \label{eq:or-m-ub}
\end{align}
Together, \eqref{eq:or-m-lb}--\eqref{eq:or-m-ub} enforce
$\widecheck{X}_{cij} = 1 \iff \widetilde{X}_{cij} \geq 1$.

\medskip
\noindent\textbf{Rule 3 (At-most-one aggregation).}
When at most one element of $\mathcal{G}_{cij}$ may equal~1, a new binary
variable $\widehat{X}_{cij}$ is introduced and linked to the sum aggregate
via:
\begin{align}
  \widehat{X}_{cij} &\leq \widetilde{X}_{cij},
    \label{eq:at-most-one-m-lb} \\
  \widetilde{X}_{cij} &\leq \widehat{X}_{cij}.
    \label{eq:at-most-one-m-ub}
\end{align}
Together, \eqref{eq:at-most-one-m-lb}--\eqref{eq:at-most-one-m-ub} enforce
$\widehat{X}_{cij} = 1 \iff \widetilde{X}_{cij} = 1$, i.e.\ if the new
variable is active, exactly one element of $\mathcal{G}_{cij}$ equals~1.

\begin{figure}[H]
\centering
\resizebox{0.6\linewidth}{!}{%
  \begin{tikzpicture}[
  every node/.style = {font=\small},
  idx/.style = {
    draw, rounded corners=3pt,
    minimum width=0.85cm, minimum height=0.60cm,
    inner sep=3pt, align=center
  },
  kept/.style  = {idx, fill=blue!12,  draw=blue!50!black},
  drop/.style  = {idx, fill=red!12,   draw=red!60!black, text=red!70!black},
  lbl/.style   = {font=\footnotesize\itshape, text=black!55},
  arr/.style   = {-{Stealth[length=6pt]}, thick, black!60},
  grpbox/.style = {
    draw=orange!70!black, rounded corners=5pt,
    fill=orange!6, inner sep=10pt, dashed, thick
  },
  derivbox/.style = {
    draw=teal!60!black, rounded corners=5pt,
    fill=teal!8, inner sep=8pt, thick
  },
]

%% ---------------------------------------------------------------
%% Step 0: root variable label at top
%% ---------------------------------------------------------------
\node[font=\normalsize\bfseries] (title) at (2.5, 0.5)
  {Index suppression};

%% ---------------------------------------------------------------
%% Root variable row (shown once, at top)
%% ---------------------------------------------------------------
\node[lbl] (rlbl) at (2.5, -0.2) {root variable $\phi$};

\node[font=\normalsize]          at (0.0, -1.0) {$\phi$};
\node[kept] (ra) at (0.8,  -1.0) {$i_1$};
\node[kept] (rb) at (1.75, -1.0) {$\cdots$};
\node[kept] (rc) at (2.7,  -1.0) {$i_k$};
\node[drop] (rd) at (3.7,  -1.0) {$j_1$};
\node[drop] (re) at (4.65, -1.0) {$\cdots$};
\node[drop] (rf) at (5.6,  -1.0) {$j_l$};

% braces
\draw[decorate, decoration={brace, amplitude=4pt},
      blue!50!black, thick]
  ($(ra.north west)+(0,0.06)$) -- ($(rc.north east)+(0,0.06)$)
  node[midway, above=5pt, font=\footnotesize, text=blue!60!black]
    {retained};

\draw[decorate, decoration={brace, amplitude=4pt},
      red!60!black, thick]
  ($(rd.north west)+(0,0.06)$) -- ($(rf.north east)+(0,0.06)$)
  node[midway, above=5pt, font=\footnotesize, text=red!60!black]
    {suppressed};

%% ---------------------------------------------------------------
%% Arrow down: "for each fixed (i1,...,ik)"
%% ---------------------------------------------------------------
\draw[arr] (2.8,-1.45) -- (2.8,-2.15)
  node[midway, right=5pt, lbl] {fix retained values $(i_1,\ldots,i_k)$};

%% ---------------------------------------------------------------
%% Aggregation group G: list of root instances varying j's
%% ---------------------------------------------------------------

% Four representative instances inside the group box
\node[font=\small] (g1) at (1.4, -2.9)
  {$\phi_{i_1,\ldots,i_k,\;j_1^{(1)},\ldots,j_l^{(1)}}$};
\node[font=\small] (g2) at (1.4, -3.55)
  {$\phi_{i_1,\ldots,i_k,\;j_1^{(2)},\ldots,j_l^{(2)}}$};
\node[font=\small] (g3) at (1.4, -4.2)
  {$\vdots$};
\node[font=\small] (g4) at (1.4, -4.7)
  {$\phi_{i_1,\ldots,i_k,\;j_1^{(s)},\ldots,j_l^{(s)}}$};

% Fit a box around them
\begin{pgfonlayer}{background}
  \node[grpbox, fit=(g1)(g2)(g3)(g4)] (gbox) {};
\end{pgfonlayer}

% Label G below the box
\node[font=\normalsize, text=orange!80!black, below=4pt of gbox]
  {$\mathcal{G}_{i_1,\ldots,i_k}$};

%% ---------------------------------------------------------------
%% Arrow right: aggregation rule
%% ---------------------------------------------------------------
\draw[arr] (gbox.east) -- ++(2.2, 0)
  node[midway, above=4pt, lbl] {aggregation rule}
  node[right] (deriv-anchor) {};

%% ---------------------------------------------------------------
%% Derived variable box
%% ---------------------------------------------------------------
\node[derivbox, right=2.2cm of gbox] (dbox)
  {$\phi_{i_1,\ldots,i_k}$};

\node[font=\footnotesize\itshape, text=teal!60!black, below=4pt of dbox]
  {derived variable};

%% ---------------------------------------------------------------
%% Annotation: "all j-tuples in J1 x ... x Jl"
%% ---------------------------------------------------------------
\node[lbl, below=0.55cm of gbox, yshift=-0.15cm]
  {$j_n^{(\cdot)} \in J_n,\quad n=1,\ldots,l$};

\end{tikzpicture}
}
\caption{Index suppression: the aggregation group $\mathcal{G}_{i_1,\ldots,i_k}$
  collects all root variable instances that share the same retained index
  values. An aggregation rule then reduces the group to a single derived
  variable.}
\label{fig:index-suppression}
\end{figure}

\end{convention}
